\section{Methods} % of the HJI Reachability Problem}
\label{sec:methods}
%
We first transform the nonlinear hyperbolic ``smoothed" HJ PDE \eqref{eq:ivp-viscous} to a linearized form and exact its unique solution based on a deduction from the solution to the heat equation. We then prescribe a sampling-based scheme, based on proximity theory, for its spatial gradients. This is found in a stochastic gradient descent algorithm in a zeroth-order optimization HJ scheme. The latter parts of this section extend the results to~\eqref{eq:HJ-RCBRT} and \eqref{eq:HJI-RCBRAT}.

\subsection{HJ Linearization by a Cole-Hopf Transformation} % of the HJI Reachability Problem}

Here, we first transform the smooth HJ equation \eqref{eq:ivp-viscous} into a linear form via the Cole-Hopf trick~\citet[\S 4.4.1]{EvansPDEBook}. The Cole-Hopf transformation that linearizes \eqref{eq:ivp-viscous} can be written as (see the derivation in Appendix \ref{app:proof_sol_hj_ivp}) 
%
\begin{align}
	\colehopf^\delta = \exp(-2\valuefunc^\delta/\delta).
	\label{eq:cole-hopf}
	\tag{Cole-Hopf}
\end{align}
%
Applying the Cole-Hopf transformation to  \eqref{eq:ivp-viscous}, we find that the fundanmental solution to the smoothed version of the Hamilton-Jacobi equation \eqref{eq:ivp-viscous} is given by \eqref{eq:ivp-viscous-solution}.
%
\begin{lemma}[Smoothed HJ Solution]
	%
	The unique solution to the smoothed version of the Hamilton-Jacobi equation \eqref{eq:ivp-viscous} is
	\begin{align}
		\label{eq:ivp-viscous-solution}
		\valuefunc^\delta(t; \state) &= -\frac{\delta}{2} \log \bigg\{(2\pi\delta t)^{-n/2}\int_{\ren} \exp \left(-\frac{1}{2\delta t}\cdot |\state-\bm{y}|^2\right) \exp \left(-\frac{2}{\delta}\cdot \valuefunc^\delta(\bm{y})\right) d\bm{y} \bigg\}
	\end{align}
	%
	\text{ in } $\ren \times (0, T]$ and $\valuefunc^\delta(t; \state) = 0$ otherwise.
	\label{lemma:hj_solution}
\end{lemma}
%
%The fundamental solution, $\fundamentalsol_{\delta t}$ to \eqref{eq:Cole-Hopf} at time $t$ is~\citet{EvansPDEBook}
%%
%\begin{align}
%\fundamentalsol_{\delta t}(\state) = \begin{cases}
%	(2 \pi \delta t )^{n/2} & \exp\left(-\dfrac{|\state|^2}{2 \delta t}\right) \text{ in } \ren \times (0, T], \\
%	\qquad 0 & \text{ otherwise }.
%\end{cases}
%\nonumber
%\end{align} 
%%
\begin{corollary}[HJ solution as a log expectation]
	On the domain of the differential equation, the unique solution can be written as the following expectation
	\begin{align}
		\valuefunc^\delta(t; \state)&= -\frac{\delta}{2} \log \bb{E}_{\statey \sim\mc{N}(\state, \delta t)}\left[\exp\left(-\dfrac{2}{\delta}\valuefunc^\delta(\bm{y})\right)\right].
		\label{eq:ivp-viscous-expectation}
	\end{align}
	%
	where $\statey$ is drawn from a Gaussian distribution with mean $\state$ and variance $\delta t$.
	\label{cor:expectation}
\end{corollary}
%
%The heat equation \eqref{eq:cole-hopf} admits the following explicit formulae (as a convolution, an integral expression, and a Gaussian expectation; see derivation in \citet[Appendix B]{HJMAD}),
%%%
%The equation \eqref{eq:cole-hopf} transforms solutions to the heat equation~\citet{EvansPDEBook, HJProx} to those of the viscous HJ PDE. It has enjoyed use in studying the energy landscape of neural networks~\citet{BaldassiLocalEntropy, Chaudhari2018} and viscous HJ loss functions~\citet{HJMAD}. 
%
Inspecting \eqref{eq:cole-hopf}, we may write
%
\begin{subequations}
\begin{align}
	\valuefunc^\delta(t; \state) = -\dfrac{\delta}{2} \log \left(\colehopf^\delta(t; \state)\right)(\state) \text{ in } \ren \times (0, T],
\end{align}  
\label{eq:value_as_colehopf}
\end{subequations}
%
from which we can write
%
\begin{align}
\partial_{\state}\valuefunc^\delta(t; \state) = - \dfrac{\delta}{2}\cdot \dfrac{	\nabla \colehopf^\delta(t; \state)}{\colehopf^\delta(t; \state)}.
\label{eq:spatial_cole_hopf}
\end{align}
%
Similarly,  for $\valuefunc_t$, recall that $\colehopf_t^\delta = \openset^\prime(\valuefunc^\delta) \valuefunc^\delta_t$, where $\openset^\prime = -(2/\delta) \exp\left(-\frac{2}{\delta}\valuefunc^\delta\right)$ is the derivative obtained from the linear transformation \eqref{eq:linear_trans}. It follows that 
%
\begin{align}
	\valuefunc_t^\delta &= \dfrac{\colehopf_t^\delta}{\openset^\prime (\valuefunc^\delta)} := \dfrac{\frac{\delta}{2} \Delta \colehopf^\delta}{-(2/\delta) \exp\left(-\frac{2}{\delta}\valuefunc^\delta\right)}
\end{align}
%
from which we can write the time derivative as 
%	
\begin{align}
	\valuefunc_t^\delta = - \dfrac{\delta^2}{4} \dfrac{\Delta \colehopf^\delta}{\exp\left(-(2/\delta) \valuefunc^\delta \right)}.
\end{align}


%Chaudhari et al.~\citet{Chaudhari2018} called $\localentropy$ the local entropy and postulated a stochastic differential equation (SDE) scheme that leverages $\localentropy$ in SGD contexts in resolving the viscous  HJ PDE. Recent publications treat HJ PDE optimization~\citet{HJProx} in light of this local entropy function %with proximal operators and Moreau envelopes provideby combining Moreau envelopes~\citet{Moreau} and proximal operator~\citet{BauschkeCombettes} algorithmics to establish zero-order optimization schemes for the viscous HJ loss function (with guaranteed convergence to global optima).


We now cast the computation of the optimal ``controllable"\footnote{We use this term because the agents in the differential game system have control inputs} reachable set into a form that leverages the Moreau envelope \eqref{eq:Moreau} in a successive approximation scheme so that the reachable set can be obtained via the subdifferential operator of the problem.

\subsection{Operators and Envelopes of HJ Equations}
%\noindent \textbf{Problem Statement} In a reachability context, we are interested in 
Monotone operator theory has had a huge impact on partial differential equations and nonlinear systems~\citet{Zarantonello71, BotAndHendrich2014}, reducing nonlinear analyses to problems of the form
%
\begin{align}
\text{ determine } \state \in \text{ zer } \bm{Q} = \{\state \in \hilbert \, | \, 0 \in \bm{Q} \state \}.
\end{align}
%
%\begin{definition}%[The Subdifferential Operator]
	The subdifferential operator~(see~\citet[\S 16]{BauschkeCombettes}), $\subdifferential: \hilbert \rightarrow 2^{\hilbert}$, of $\bm{g}(\state)$ at $\bar{\state}$ is the set of all $\subdiffset$ that satisfies
%
\begin{align}
\{\subdiffset\in \hilbert \mid (\forall \state \in \hilbert) \langle \state - \bar{\state} , \subdiffset\rangle + \term(\bar{\state}) \le \term(\state)\}.
\label{eq:subdifferential}
\end{align}
%
That is, $\term$ is subdifferentiable at $\state$ if $\subdifferential(\state) \neq \emptyset$; hence, elements of $\subdifferential(\state)$ are the subgradients of $\term$ at $\state$.
%\end{definition}

An optimization implication of \eqref{eq:subdifferential} is Fermat's rule, which states that for $\term: \hilbert \rightarrow (0, T)$~\citet{Moreau65},
%
\begin{align}
\argmin \,\, \term = \zeroes \subdifferential.
\end{align}
%
The proximity operator of $\term(\state)$ at a time, $t \in (0, T]$ is the minimizers set~\citet{MoreauDualCvxProxPts} 
%
\begin{align}
\prox_{t\term}(\state) \triangleq \arg \min_{\bm{z} \in \hilbert }  \left(\term(\bm{z}) + \frac{1}{2t}\| \bm{z} - \state \|^2\right)
\tag{Prox-Op}
\label{eq:proximal}
\end{align}
%
that define the Moreau envelope $\valuefunc(\cdot, t)$, defined as~\citet{Moreau}
%
\begin{align}
\valuefunc(\state, t) = \min_{\bm{z} \in \hilbert }  \left(\term(\bm{z}) + \frac{1}{2t}\| \bm{z} - \state\|^2\right).
\tag{Moreau-Env}
\label{eq:Moreau}
\end{align}
%

The envelope $\valuefunc$ expands troughs of $\term(\state)$ and its local minimizers concur with those of $\term$~\citet{Chaudhari2018}. 
%
\begin{lemma}
	Suppose that $\term(\state)$ is BUC, its subdifferential exists, the set $\mc{S}_\alpha \triangleq \{z \in \ren: \term(\bm{z}) \le \inf \term + \alpha \}$ is compact, $\zeroes \subdifferential$ exists and $\state \in \mc{S}_\alpha$. An important result established in~\citet[Lemma A.3]{HJMAD} relates the spatial gradients of the solution to \eqref{eq:ivp} to the proximity operator of $\term(\state)$ at time $t$ as  
	%
	%\begin{subequations}
		\begin{align}
		\prox_{t\term}(\state) &= \state - t \partial_{\state}  \valuefunc(t; \state) :\approx \state -  t \cdot \partial_x\valuefunc^\delta(t; \state).
		%\label{eq:prox-hj-approx}
		\label{eq:prox-hj}
		\end{align}
	%\end{subequations}
\end{lemma}
%Our goal now is to leverage the Proximity operator in the minimization \eqref{eq:distance} while playing the zero-sum two-player differential game \eqref{eq:HJI-RCBRAT}  to optimize for the RCBRT~\eqref{eq:rcbrt} and the zero levelset~\eqref{eq:reachables}. In this light, we can replace $\valuefunc(\state, t)$~\eqref{eq:ivp} by its viscous and smooth proxy~\eqref{eq:ivp-viscous}, evaluate the minimizer \eqref{eq:distance} using the proximal operator result~\eqref{eq:prox-hj} and solve for the \texttt{inf-sup} optimal control problem in \eqref{eq:value_lower}. %  $\valuefunc(\state, t)$. 
%Let us state the following preliminary result in the build up to our results.
%%
%\begin{lemma}
%	The minimizers of $\term(\state)$ are the zeros of its subdifferential \ie $\Argmin \term(x) = \zeroes \subdifferential :\triangleq \{\state \in \hilbert \mid 0 \in \subdifferential (x)\}$. 
%\end{lemma}
%%
%The above is a statement of Fermat's Rule~\citet[Theorem 16.1]{BauschkeCombettes}.
%
%An intertwined relation between the subdifferential operator and the monotone proximal operator~\eqref{eq:proximal} is now stated.
%\begin{proposition}
% For all $ (\state, p) \in \hilbert$  %~\citet{CombettesMonoOpTheory}, \ie 
%	%
%	\begin{align}
%	 \prox_{tf}(\state + \subdiffset) 	\Leftrightarrow  (\forall \state \in \hilbert) \,	(\forall \subdiffset \in \hilbert), \,\,\subdiffset \in \subdifferential
%	\end{align}
%	%
%	or
%	\begin{align}
%	p \in \prox_{t\term}(\state) \Leftrightarrow 	%\text{prox}_{tf}(\state - p) 
%	(\state - p) \in \subdifferential(p). 
%	\end{align}
%	%
%	Proposition \ref{prop:prox_sub} is a restatement of Proposition 16.34 in~\citet{BauschkeCombettes}. 
%	This will enable us to reformulate the HJI equality~\eqref{eq:Reach-HJ} as a successive approximation problem when finding the extremas of $\valuefunc(\state, t)$ via optimization  for
%	%
%	\begin{align}
%	\state_0 \in \hilbert \text{ such that  } \forall\, k \in \bb{N}, \,\, \state_{k+1} = \prox_{t\term}(\state_k).
%	\end{align}
%	\label{prop:prox_sub}
%\end{proposition}
%%
%%This will be shown to allow us solve a structured optimization problem using proximal operators of the functions present in the HJI equality \eqref{eq:Reach-HJ}.
%\begin{proposition}[Prop 16.21~\citet{CombettesMonoOpTheory}]
%	For a BUC $\term(\state)$, a fixed point of the proximity values of $\term(\state)$ is equivalent to the zeroes of the proximity operator of $\term(\state)$, \ie 
%	%
%	\begin{align}
%		\Argmin \term(x) = \zeroes \subdifferential = \subdifferential^\star(0)  = \text{Fix} \prox_{t\term} = \zeroes \prox_{t\term^\star}.
%	\end{align}
%\end{proposition}

\subsection{HJ PDE's spatial gradients and its proximity operator}
%
Consider the kernel $\colehopf^\delta(\state, t)$ in \eqref{eq:Cole-Hopf} which corresponds to the solution of \eqref{eq:ivp-viscous}. Resolving the convolution \eqref{eq:heat_sol} is intractable for large $n$ and time step $t$. However, it is same as the probability density of a Gaussian distribution with mean $\state$ and variance $\delta t$ (as derived in~\citet[Appendix B]{HJMAD}), \ie
%
\begin{align}
\partial \colehopf_{\state}^\delta(t; \state) = -\dfrac{1}{\delta t} \cdot \bb{E}_{y\sim\mc{N}(\state, \delta t)} \left[(\state - \bm{y})\,\exp(\term(\bm{y})/\delta)\right].
\label{eq:heat_gaussian}
\end{align}
%
Plugging \eqref{eq:heat_sol_expect} and \eqref{eq:heat_gaussian} into \eqref{eq:spatial_cole_hopf}, we find that 
%
\begin{align}
t \cdot \partial \valuefunc_{\state}^\delta(t; \state) %&=  \dfrac{\bb{E}_{y\sim\mc{N}(\state, \delta t)} \left[(\state - \bm{y})\,\exp(\term(\bm{y})/\delta)\right]}{\bb{E}_{y\sim\mc{N}(\state, \delta t)}\left[\exp(-\delta^{-1} \, \term(\bm{y}))\right]}  \\
%
&=  \state -\dfrac{\bb{E}_{y\sim\mc{N}(\state, \delta t)} \left[ \bm{y}\,\exp(-\term(\bm{y})/\delta)\right]}{\bb{E}_{y\sim\mc{N}(\state, \delta t)}\left[\exp(-\delta^{-1} \, \term(\bm{y}))\right]}
\label{eq:value_grad_expect}
\end{align}
%
provides the time-scaled spatial gradients of $\valuefunc(t; \state)$ in \eqref{eq:hji_prox_pre} without the need for unstableschemes employed in upwinding methods within grid discretization numerical reachable sets' computational algorithms~\citet{levpyconf, LevelSetTOMS, MitchellLSToolbox}.

Substituting the r.h.s of \eqref{eq:value_grad_expect} into \eqref{eq:prox-hj}, we find that 
%%
%\begin{align}
%t \cdot \nabla_{\state}\valuefunc^\delta(t; \state) = \state - \prox_{t\term}(\state)
%\label{eq:prox-hj-rearranged}
%\end{align}
%
%which upon substitution into \eqref{eq:value_grad_expect} implies that 
%
\tcb{
	\begin{align}
	%t \cdot \partial \valuefunc_{\state}^\delta(t; \state) 
	\prox_{t\term}(\state) &=   \dfrac{\bb{E}_{y\sim\mc{N}(\state, \delta t)} \left[\bm{y} \cdot \exp(-\term(\bm{y})/\delta)\right]}{\bb{E}_{y\sim\mc{N}(\state, \delta t)}\left[\exp(-\delta^{-1} \, \term(\bm{y}))\right]} %\nonumber \\
	%
	%&:= \prox_{t\term}(\state).
	\label{eq:prox_value_grad}
	\end{align}
}{Proximal Operator for $\partial_{\state} \valuefunc^\delta(t; \state)$}


\subsection{Hamilton-Jacobi--Proximity Operator Relation}
%
Now consider the HJI PDEs~\eqref{eq:HJ-RCBRT} and ~\eqref{eq:HJI-RCBRAT} in light of \eqref{eq:prox_value_grad}. We can write the domain of these PDEs in light of the proximal expression \eqref{eq:prox_value_grad} as 
%
\begin{align}
&\partial \valuefunc_t(t; \state) + \min\{0, \hamfunc\left(x, \partial \valuefunc_{\state}(t; \state)\right)\} = 0, \nonumber \\
%\valuefunc(0; \state) &= \bm{g}(0; \state) \\
%
&\approx \partial \valuefunc^\delta_t(t; \state) + \min\{0, %
\hamfunc^\delta\left(t; x, \partial \valuefunc^\delta_{\state}\right)\} = 0, \nonumber \\
%
&:= \partial \valuefunc^\delta_t(t; \state) + \min\bigg\{0, %
\max_{\control \in \mathcal{U}} \min_{\disturb \in \mathcal{W}} \, \bigg\langle f(t; \state, \bm{u}, \bm{w}),  \frac{1}{t}(\state - \prox_{t\term}(\state)) \bigg\rangle
\bigg\} = 0,
\end{align}
%
for \eqref{eq:HJ-RCBRT}.
%
Similarly, the right-hand-side of the variational inequality for \eqref{eq:HJI-RCBRAT} can be written as %in terms of \eqref{eq:prox-hj-rearranged} as 
%
%\begin{subequations}
	\begin{align}
	&\min\{\partial \valuefunc_t(t; \state) + \hamfunc\left(t; x, \partial \valuefunc_{\state}\right), \bm{g}(\state, t) - \valuefunc(t; \state)\}, \nonumber \\
	%
	&\approx \min\{\partial \valuefunc^\delta_t(t; \state) + \hamfunc^\delta\left(t; x, \partial \valuefunc^\delta_{\state}\right), \bm{g}(\state, t) - \valuefunc^\delta(t; \state)\}, \nonumber \\
	%
	%&:= \min\{\valuefunc^\delta_t(t; \state) + \hamfunc^\delta\left(t; x, \partial \valuefunc^\delta_{\state}\right), \bm{g}(\state, t) - \valuefunc^\delta(t; \state)\} \nonumber \\
	%
	&:= \min\{\partial \valuefunc^\delta_t(t; \state) +  \max_{\control \in \mathcal{U}} \min_{\disturb \in \mathcal{W}} \, \langle f(t; \state, \bm{u}, \bm{w}), \partial \valuefunc_{\state}^\delta  \rangle,  \bm{g}(t; \state) - \valuefunc^\delta(t; \state)\}. 
	\label{eq:hji_prox_pre}
	\end{align} 
%\end{subequations}
%
Thus, evaluating the spatial gradient of $\valuefunc^\delta$ in \eqref{eq:hji_prox_pre} amounts to drawing samples using the proximity expression above and using \eqref{eq:heat_sol_expect} to evaluate $\valuefunc^\delta$. We can now write  \eqref{eq:HJI-RCBRAT}  as
%
%\begin{subequations}
	\begin{align}
	\min&\bigg\{\partial \valuefunc^\delta_t(t; \state) + \max_{\control \in \mathcal{U}} \min_{\disturb \in \mathcal{W}}\, \bigg\langle f(t; \state, \bm{u}, \bm{w}), \frac{1}{t}(\state - \prox_{t\term}(\state)) \bigg\rangle,  \\ %(1/t)\prox_{t\term}(\state) \rangle, \nonumber \\
	\bm{g}(\state, t) &+ \delta \log \bb{E}_{\bm{y}\sim \mc{N}(\state, \delta t)}\left[\exp(-\delta^{-1} \term(\bm{y}))\right]\bigg\} \le 0, \quad \valuefunc(0; \state) &= \term(0; \state).
	\end{align}
%\end{subequations}
%%
%\begin{subequations}
%	\begin{align}
%	\min&\bigg\{\partial \valuefunc^\delta_t(t; \state) + \max_{\control \in \mathcal{U}} \min_{\disturb \in \mathcal{W}}\, \bigg\langle f(t; \state, \bm{u}, \bm{w}), \nonumber \\
%	%
%	& \qquad\quad  \dfrac{\bb{E}_{y\sim\mc{N}(\state, \delta t)} \left[\bm{y} \cdot \exp(-\term(\bm{y})/\delta)\right]}{\bb{E}_{y\sim\mc{N}(\state, \delta t)}\left[\exp(-\delta^{-1} \, \term(\bm{y}))\right]} \bigg\rangle,  \\ %(1/t)\prox_{t\term}(\state) \rangle, \nonumber \\
%	 \bm{g}(\state, t) &+ \delta \log \bb{E}_{\bm{y}\sim \mc{N}(\state, \delta t)}\left[\exp(-\delta^{-1} \term(\bm{y}))\right]\bigg\} \le 0 \nonumber \\
%	%
%	\valuefunc(0; \state) &= \term(0; \state).
%	\end{align}
%\end{subequations}

\subsection{Problem Algorithmics} % samples for Robust Reachable Sets Computation}

%
\begin{savenotes}
	\begin{algorithm}[t!]
		\caption{Sampling-based RCBRT descent.\label{alg:rcbrt-sample}}
		\begin{algorithmic}
			\State Given $\{\epsilon, \delta\}>0$, and initial step size $\alpha$.
			\Function{\texttt{\textbf{ProxReach}}}{$\state, \term, \eta, t_0, T$} \Comment{Sample size, $\eta$.}
			%\State $\quad$ \Comment{value function, and interval $(t_0, T]$.}
			\State \texttt{Set $\tau= t_0, \, t_s = (T-t_0)/k$}.
			%
			\While{	$T - \tau > \epsilon T$}
			%\For{$i \in \bb{N}$}
			\State $\bm{y}[\tau] \leftarrow \mc{N}(\state, \delta)$  \label{alg:rcbrt::sample}; \Comment{Draw finite samples from $\state$.}
			%
			\State $\bm{z}[\tau] \leftarrow \term(\bm{y}[\tau])$; \Comment{See equation \eqref{eq:value_as_colehopf}}.
			%
			\State $\partial_{\bm{z}}\valuefunc(\tau; \bm{z}) \leftarrow \texttt{Prox}_{\tau\term}(\bm{z}, \delta)$ \Comment{Eq. \eqref{eq:proximal}.} \label{alg:rcbrt:prox_update}
			\State $\tau \leftarrow \tau + t_s$
			\State $\statez[\tau + t_s] \leftarrow \statez[\tau] - \alpha \cdot \term[\tau]$; \label{alg:rcbrt::time-adjuster}	
			\State $\tau \leftarrow \texttt{Adjuster}	(\tau, \term[\tau], \term[\tau-\Delta t])$;
			%\EndFor
			%	\State $\prox_{t\term}(\state), t \leftarrow \text{ MLP}(-\bm{z}/\delta, [\bm{y}^1,\cdots, \bm{y}^N])^\top$	
			\EndWhile
			\EndFunction
		\end{algorithmic}
	\end{algorithm}
\end{savenotes}
We now cast the solution of the variational HJI computational problem as a Monte Carlo algorithmics. Towards finding global solutions to the optimization problem, we adapt the Hamilton-Jacobi Moreau adaptive descent algorithm of~\citet{HJMAD} to the computation of the HJI\eqref{eq:HJ-RCBRT} and \eqref{eq:HJI-RCBRAT}. %Looking at equation \eqref{eq:prox_value_grad}, since we are working in a sampling-based values spatial derivative setting, we must avoid using too large samples to mitigate against time inefficiency while solving the problem. 

The Moreau adaptive descent algorithm works by minimizing the Moreau envelope~\eqref{eq:Moreau} of the HJ PDE in a nonconvex optimization scheme via a zeroth-order sampling scheme. In its native form, within a set convergence loop, it works by 
%
\begin{itemize}
	\item first computing the spatial gradient of $\valuefunc^\delta(t; \state)$; then
	%
	\item updating the states of the dynamic system in a gradient descent scheme (under an appropriate step size); and lastly
	%
	\item uses an adaptive forward-backward time-update scheme to update the index of time.
	%
\end{itemize}
%

The algorithm is succintly described in Algorithm \autoref{alg:rcbrt-sample}. Given a large initial step size $\alpha$, time interval $(t_0, T]$ and the state values for a given froblem $\state$, the algorithm proceeds in a reactive sampling-based gradient descent fashion by drawing finite samples of size $\eta$ from the (potentially large) state space $\state \in \ren$. The samples on Line \ref{alg:rcbrt::sample} of Algorithm \autoref{alg:rcbrt-sample} are draw from a multinomial Gaussian filter for a total of $\eta$ samples whose  mean and variance are $\state$ and $\delta t$, respectively in every iteration of the algorithm. Within line \ref{alg:rcbrt:prox_update} of Algorithm \autoref{alg:rcbrt-sample}, the Hamiltonian for the system is computed but negated to make the computational scheme conform tot he backward reachable setting (since the iterative loop proceeds forward in the index of time in Algorithm \autoref{alg:rcbrt-sample}). In all our implementations, the time index within the loop is adaptively adjusted (\ie Line \ref{alg:rcbrt::time-adjuster}) based on a version of the Armijo-Goldstein condition~\citet{Armijo1966minimization} to ensure that an appropriate time step is taken in every iteration as we seek the minimum value of the Moreau envelope.

After convergence, the solution(s) $\{\state\}$ that minimize the value function are returned. We now adaopt the HJ-MAD algorithm to the resolution of HJ problems with adversarial control inputs in a two-player differential game. This is delineated in Algorithm~\ref{alg:rcbrt-sample}.
%While a weighted moving average of the gradient was mentioned in~\citet{HJMAD}, no incorporation of prior data was imposed on the computational scheme. In this work, we incorporate a data-driven prior to every successive sample by first fitting Gaussian models to the numerator and denominator in \eqref{eq:prox_value_grad}. Preventing large samples being used to compute the proximal operator then amounts to conditioning the Gaussian model as $(\prox_{(t+1)\term}|\prox_{t\term})$ at each time step $t$. This is similar to linear regression with the advantage of allowing us to integrate prior information in the form of a Gaussian-inverse-Wishart distribution model. 
After the optimization of the extrema of the value function, finding the safe set is equivalent to executing a search algorithm within the optimized value function to evaluate where the condition \eqref{eq:reachables} are fulfilled. In the next section, we will connect the scheme we have presented so far to the safe reachability resolution problem of two rockets on a cartesian grid as originally used in our previous paper~\citet{levpyconf}.
%For an \textit{equivalent sample size} $\kappa_0$ and its belief, $\nu_0$ on the prior variance, $\Sigma_0$ that is associated with the precision $\Lambda_0 = \Sigma_0^{-1}$ of a dataset $D$, suppose that the mean and variance of the normal-inverse-Wishart natural conjugate prior~\citet{WishartPrior} are, 
%%
%\begin{align}
%	\Sigma \sim IW_{\nu_0}(\Lambda_0^{-1}), \quad \mu\mid \Sigma \sim \mc{N}(\mu_0, \Sigma/\kappa_0)
%\end{align} 
%%
%with associated empirical estimates $\hat{\Sigma}$ and $\hat{\mu}$. Then, the natural conjugate prior of the normal-inverse-Wishart is given as $	p(\mu, \Sigma) \triangleq NIW(\mu_0, \kappa_0, \Lambda_0, \nu_0)$.
%%
%%\begin{subequations}
%%	\begin{align}
%%		p(\mu, \Sigma) \triangleq NIW(\mu_0, \kappa_0, \Lambda_0, \nu_0).
%%	\end{align}
%%\end{subequations}
%We are concerned with its maximum-a-posteriori 
%%
%\begin{align}
%	p(\mu,& \Sigma | D, \mu_0, \kappa_0, \Lambda_0, \nu_0) =  NIW(\mu, \Sigma | \mu_n, \kappa_n, \Lambda_n, \nu_n), \nonumber \\
%	 \kappa_n &= \kappa_0 + n, \,\, \nu_n = \nu_0 + n  
%\end{align}
%%
%where 
%\begin{subequations}
%	\begin{align}
%		\mu_n &= \dfrac{\kappa_0\mu_0}{\kappa_0 + n} + \dfrac{n}{\kappa_0 + n}\hat{\mu}, \\
%		%
%		\Lambda_n &= \Lambda_0 +S +\dfrac{ \kappa_0 n (\kappa_0 + n) (\hat{\mu} - \mu_0)(\hat{\mu}-\mu_0)^\top}{\kappa_0 + n}, \\
%		%
%		S &= \sum_{i=1}^{n} (\state - \bar{\state})(\state - \bar{\state})^\top.
%	\end{align}
%	\label{eq:normal_inv_wishart}
%\end{subequations}
% %Suppose that we set $\hat{\Lambda}t$ as the empirical variance of the computed proximity operator, and $\hat{\mu}$ as the empirical mean. For prior parameters, $\Phi, \mu_0, n_0$, and $m$, the normal-inverse Wishart prior
%
%We recursively incorporate new samples by leveraging the MAP parameters given in \eqref{eq:normal_inv_wishart} to new samples drawn from \eqref{eq:prox_value_grad}. 

\subsection{Convergence analysis of the algorithm}
\todo{Haoxiang, please follow the results in this link: }